\begin{abstract}
当点光源绕物体移动时，物体外观的许多变化取决于光照落在何处。
现有可重光照的高斯泼溅方法逐个图元估计这一过程；我们则在光源的成像平面上进行建模。
从光源视角光栅化高斯时，alpha 合成权重恰好表示各图元截获的光源像素通量比例。
\ours 将这一沉积量存入由深度矩和可学习通量特征组成的光源空间图集。
任意接收点——训练时的高斯或渲染时的像素——均可读取图集金字塔，得到经小型网络修正的滤波阴影检验，以及对沉积特征呈线性的光传输项。
此外，在几何形成阶段，训练需要在线性辐亮度域而非显示编码域计算损失；在两个包含细小、密集结构的场景上，受控实验表明这一选择使留出视图的 PSNR 提高了 6--7\,dB。
在三个公开基准的 18 个场景上，\ours 的平均 PSNR 达到 33.47\,dB，高于四种近期基线中最强方法的 31.59\,dB；它也是唯一在每个基准上均位列前二的方法，训练耗时约 20 分钟。
消融实验同时界定了光传输项的作用范围：它对半透明材质最为重要，而在多数其他材质上，相同规模的局部网络即可达到相当效果。
\end{abstract}
